\begin{definition}
    Базисом в линейном пространстве $L$ называется конечная, упорядоченная, ЛНЗ система векторов из $L$ такая, что любой вектор из $L$ по ней раскладывается.
\end{definition}
Из критерия ЛНЗ разложение любого вектора по базису единственно. Координатами вектора будем называть коэффициенты в его разложении по базису. По договоренности координаты векторов записываются в столбец. Формы записи:
\[
(\vec{e_1}, \dots, \vec{e_n}) \begin{pmatrix} x_1 \\ \vdots \\ x_n \end{pmatrix}=x_1 \vec{e_1} + \dots +x_n \vec{e_n}=ex
\]
Заметим, что в нулевом пространстве нет базиса, так как базис обязан быть непустым.
\begin{proposition}
    При фиксированном базисе в линейном пространстве $L$:
    \[
    \forall \alpha \in \R: \forall\vec{x}, \vec{y} \in L \hookrightarrow \alpha\vec{x}=\begin{pmatrix}
        \alpha x_1 \\
        \vdots \\
        \alpha x_n
    \end{pmatrix}, \vec{x} + \vec{y}=\begin{pmatrix}
        x_1 + y_1 \\
        \vdots \\
        x_n + y_n
    \end{pmatrix}
    \]
\end{proposition}
\begin{proof}
    Следует из ассоциативности умножения на скаляр и дистрибутивности сложения: $\alpha (x_1 \vec{e_1} + \dots +x_n \vec{e_n})=(\alpha x_1) \vec{e_1} + \dots + (\alpha x_n) \vec{e_n}$ и $x_1 \vec{e_1} + \dots +x_n \vec{e_n} + y_1 \vec{e_1} + \dots + y_n \vec{e_n}= (x_1 + y_1) \vec{e_1} + \dots + (x_n + y_n) \vec{e_n}$.
\end{proof}
\begin{proposition}
    Набор векторов ЛЗ тогда и только тогда, когда набор их координатных столбцов ЛЗ.
\end{proposition}
\begin{proof}
    Из предыдущей теоремы следует, что линейная комбинация векторов равносильна линейной комбинации их координатных столбцов, отсюда получаем требуемое.
\end{proof}
\begin{theorem}
    Пусть в $L$ есть базис из $n$ векторов, тогда любой набор из $m>n$ векторов линейно зависим.
\end{theorem}
\begin{proof}
    Пусть $\vec{e_1}, \dots, \vec{e_n}$ - базис, а $\vec{f_1}, \dots, \vec{f_m}$ произвольный набор из $m>n$ векторов. Рассмотрим матрицу из координатных столбцов $\vec{f}$ в базисе $\vec{e}$: её размеры $n \times m$, а значит её $rg \leqslant n<m$, следовательно, набор $\vec{f_1}, \dots, \vec{f_m}$ - ЛЗ.
\end{proof}
\begin{corollary}
    Если существует базис из $n$ векторов, то и все другие базисы будут состоять из $n$ векторов.
\end{corollary}
\label{Размерность.}
Количество векторов в базисе называется \textit{размерностью} пространства. Обозначение: $\dim{L}=n$. Если же в $L$ для любого $m$ существует ЛНЗ набор из $m$ векторов, то $L$ называется <<бесконечномерным>>. В таком пространстве мы не определяем базис. Примеры таких пространств: пространство всех многочленов, пространство непрерывных на отрезке функций. В дальнейшем мы будем рассматривать в основном конечномерные пространства.
\begin{proposition}
    Пусть $\dim{L}=n$, тогда любой ЛНЗ набор из $n$ векторов является базисом.
\end{proposition}
\begin{proof}
    Пусть $\vec{f_1}, \dots, \vec{f_m}$ произвольный набор из $n$ ЛНЗ векторов. Тогда в силу теоремы 4.1 добавив к этому набору любой вектор $\vec{x} \in L$ мы получим ЛЗ набор, т.е. $\alpha_1 \vec{f_1} + \dots +\alpha_n \vec{f_n} + \beta \vec{x}=0$, причем $\beta^2+\sum_{i=1}^n \alpha_i^2 >0$ и $\beta \neq 0$ из ЛНЗ изначального набора. Следовательно $\vec{x}$ расскладывается в линейную комбинацию векторов $\vec{f}$: $\sum_{i=1}^n -\frac{\alpha_i}{\beta} \vec{f_i}=\vec{x}$.
\end{proof}
\begin{proposition}
    Пусть  $\dim{L}=n$, тогда любой ЛНЗ набор из $m<n$ векторов можно дополнить до базиса.
\end{proposition}
\begin{proof}
    Так как любой ЛНЗ набор из $m<n$ векторов не является базисом, то существует вектор, который не выражается через данный набор, а значит этот вектор можно добавить в набор с сохранением ЛНЗ. Продолжаем так рассуждать, пока размер набора не станет равен $n$.
\end{proof}
\subsection{Матрица перехода}
Пусть даны линейное пространство $L$, его базисы $e, e'$ (строки из базисных векторов), вектора $x, x' \in L$.
\begin{definition}
    Матрицей перехода называется матрица $S$, составленная из координатных столбцов базисных векторов из $e'$ в базисе $e$.
\end{definition}
Приведем несколько тождеств, следующих из определения:
\[
    e'S^{-1} =e, \ \ e'=eS, \ \ \ \ \ \vec{a}=ex=e'x'=eSx' \Rightarrow x=Sx' 
\]
\begin{proposition}
    Пусть $dim \ L = n$ и $e$ - базис в $L$. Тогда любая невырожденная матрица $S_{n \times n}$ является матрицей перехода от $e$ к некоторому базису $e'$.
\end{proposition}
\begin{proof}
    Так $e$ - ЛНЗ и $S$ невырожденная, то и $e'=eS$ - ЛНЗ. А так как любой вектор $\vec{a}$ выражался через $e$, то  $\vec{a}=ex=e'x'=e'S^{-1}x$.
\end{proof}
\subsection{Линейное подпространство}
Пусть $L$ - линейное пространство.
\begin{definition}
    Непустое подмножество $L'$ множества $L$ называется подпространством $L$, если выполнено: $\forall x, y \in L' \hookrightarrow x+y\in L'$ и $\forall \alpha \in \R, \forall x \in L' \hookrightarrow \alpha x \in L'$.
\end{definition}
Перечислим несколько следствий из определения:
\begin{itemize}
    \item Так как $L'$ непустое, то $\exists \vec{x}\in L \Rightarrow 0 \cdot \vec{x}=0 \in L'$. Значит $0 \in L'$.
    \item Прочие аксиомы линейного пространства выполнены и для $L'$, так как $L'\subset L$. Значит $L'$ - линейное пространство.
    \item $L$ является подпространством $L$.
    \item Нулевое пространство является подпространством любого пространства $L$.
\end{itemize}
Примеры подпространств: многочлены степени не выше 2025; множество векторов, параллельных заданной плоскости в пространстве $\R^3$.
\subsection{Линейная оболочка}
Пусть $L$ - линейное пространство и $P \subset L$ - непустое подмножество.
\begin{definition}
    Линейная оболочка множества $P$ - множество всевозможных конечных линейных комбинаций элементов $P$.
\end{definition}
Обозначим её $\tilde{L}$, иногда будем использовать обозначение $\langle P\rangle$. Из определения $\tilde{L} \subset L$, а значит $\tilde{L}$ - линейное подпространство.
\begin{theorem}
    Пусть $P \subset L$ и $\tilde{L}=\langle P\rangle$. Тогда если $\exists f_1, \dots, f_k $ - ЛНЗ $\in P:\forall n\in P \hookrightarrow n=\alpha_1 f_1+\dots+\alpha_k f_k$, то набор $f_1, \dots, f_k$ - базис в $\tilde{L}$.
\end{theorem}
\begin{proof}
    По определению набор $f_1, \dots, f_k$ - базис в $P$, а множество $\tilde{L}$ получается объединением линейных комбинаций элементов из $P$, значит любой вектор из $\tilde{L}$ выражается через набор $f_1, \dots, f_k$. Тогда по определению он является базисом и для $\tilde{L}$.
\end{proof}
Простое следствие: если $P$ - конечное множество из $k$ векторов, то $\dim{\langle P}\rangle \leqslant k$.
\begin{proposition}
    Пусть $L' \subset L$. Тогда если $\dim{L}'=\dim{L}$, то $L'=L$.
\end{proposition}
\begin{proof}
    Зафиксируем базис из $n$ векторов в $L'$. Тогда в силу утверждения 4.3 он же является базисом для $L$. Значит $L'=L$.
\end{proof}
Пусть ${e_1}, \dots, {e_k}$ - базис в $L'$ и $k<n$. Дополним его векторами ${e_{k+1}}, \dots, e_n$ до базиса в $L$. Тогда получаем критерий: \[
\vec{x} \in L' \Leftrightarrow x_{k+1}=\dots=x_n=0 \Leftrightarrow\left\{
\begin{array}{ll}
x_{k+1}=0 \\
\vdots\\
x_n=0
\end{array}
\right.
\]
Переходя к другому базису получаем другую, менее тривиальную систему. Значит мы можем задать $k$-мерное подпространство с помощью системы из $n-k$ линейных уравнений.